% !TeX root = main.tex
\section{Background}

\subsection{LLM Guidance and Semantic Grounding}

먼저 지침 생성과 실행 사이의 표현을 구분한다. $d_t$는 시점 $t$에 LLM으로 보내는
환경의 텍스트 요약이고, $G_r$는 갱신 시점 $r$에 LLM이 생성한 지침이다. 지침은
자연어 서브골과 그 우선순위 $\{(g_k,\rho_k)\}_{k=1}^{K_r}$, 그리고 선택적인 행동
규칙으로 구성된다. 이 상태에서 서브골은 아직 학습 보상으로 직접 사용할 수 없다.

이를 위해 각 환경은 유한한 접지 라이브러리
$\mathcal L=\{\phi_j\}_{j=1}^{M}$를 제공한다. 이 논문에서 \emph{predicate}는
전이 $(s_t,\mathbf a_t,s_{t+1})$의 특정 속성을 측정하는 불리언 또는 실숫값 함수
$\phi_j$를 뜻한다. 그 시점의 함수 출력은
$f_{j,t}=\phi_j(s_t,\mathbf a_t,s_{t+1})$로 표기한다. 접지 라이브러리는 실행 전에
정의된 측정 함수의 목록이며 온라인에서 LLM이 새 함수를 만들어 추가하는 구조가 아니다.
접지 모듈은 각 서브골 $g_k$를 라이브러리의 함수에 대응시키고 중복 항목을
합쳐, 활성 predicate와 비음수 가중치의 집합
$\mathcal I_r=\{(j,w_{r,j})\}$를 만든다. 따라서 의미적 셰이핑 보상은

\begin{equation}
F_t=\operatorname{clip}\!\left(
  \sum_{(j,w_{r,j})\in\mathcal I_r}w_{r,j} f_{j,t},-c,c
\right)
\label{eq:shaping}
\end{equation}

로 계산한다. 여기서 $c$는 한 전이에서 보조 신호가 지나치게 커지는 것을 막는
절댓값 상한이다. 저수준 학습 보상은

\begin{equation}
r_t^{\mathrm{train}}=r_t^{\mathrm{env}}+\lambda_tF_t
\label{eq:reward}
\end{equation}

로 구성하며, $r_t^{\mathrm{env}}$는 원래 환경 보상이고 $\lambda_t$는 보조 신호의
영향을 조절하는 계수이다. 행동 규칙은 특정 행동을 금지하는 마스크(hard mask)나
행동 선택 점수를 높이는 선호 가중치(soft preference)로 변환될 수 있다. RSVP의
주된 실험은 갱신 효과를 분리하기 위해 셰이핑 경로에 집중하며, 행동 마스킹은 별도의
제거 실험으로 다룬다.

\section{Problem Formulation and Motivation}

\subsection{Guidance Refresh as a Decision Problem}

고정 정책은 마지막 성공 호출 이후 $F$스텝마다 Commander를 갱신한다. 반면 우리가
구하려는 갱신 정책 $\eta$는 현재 상태, 보유 지침, 학습된 감시 신호를 이용해 호출
여부를 결정한다. 학습 예산 $N$에서의 목적은 개념적으로

\begin{equation}
\max_{\eta}\quad
\mathbb E[J_N^{\mathrm{env}}(\eta)]
-\alpha\,\mathbb E[C_N(\eta)]
\label{eq:objective}
\end{equation}

로 쓸 수 있다. $J_N^{\mathrm{env}}$는 셰이핑 보상이 아닌 외부 환경 보상으로 측정한
정책 성능이며, $C_N$은 실제 LLM 요청, 토큰, 지연 등의 비용이다. 식
~\eqref{eq:objective}는 연구 목표를 나타내며, 현재 구현이 이 목적함수를 직접
최적화하거나 호출 예산을 보장한다는 뜻은 아니다.

단순한 상태 변화량이나 최근 누적 보상도 갱신 신호로 사용할 수 있다. 그러나 같은 상태
변화도 목표 달성 지침과 위험 회피 지침에서 의미가 다르고, 누적 보상의 저하는 지침의
노후화뿐 아니라 탐색 실패나 정책 학습의 불안정성도 반영한다. RSVP는 보유 지침의
가중치로 감시 대상을 선택하여 이러한 내용 비의존성을 줄이고자 한다.

\subsection{Why Adaptive Refresh Matters}

첫째, 고정 주기의 선택은 단순한 시스템 파라미터가 아니라 학습 데이터의 분포를
바꾸는 결정이다. 셰이핑 지침이 바뀌면 어떤 전이에 추가 보상이 주어지는지가 달라지고,
행동 마스킹을 사용할 경우 탐색 궤적 자체도 변한다. 잘못된 갱신 주기는 비용만이 아니라
최종 정책에도 영향을 줄 수 있다.

둘째, LLM 호출 비용이 감소하더라도 불필요한 호출을 줄여야 할 이유는 남는다. 호출은
지연과 에너지뿐 아니라 확률적인 전략 재표현을 도입한다. 유효한 지침을 너무 자주
교체하면 저수준 학습기가 일관된 보조 목적을 따라갈 시간을 잃을 수 있다. 반대로 오래된
지침을 무기한 유지하면 변화한 국면과 학습 신호가 어긋날 수 있다. 필요한 것은 단순한
호출 최소화가 아니라 지침이 유효한 동안 재사용하고 변화가 누적될 때 갱신하는 것이다.

셋째, 이 문제는 특정 LLM이나 과업 도메인에만 한정되지 않는다. 느린 고수준 모듈의 출력을
빠른 저수준 제어기가 여러 스텝 재사용하는 계층형 시스템에서는 언제 계획을 다시
계산할지 결정해야 한다. 지침의 실행 의미와 연결된 감시 신호를 사용한다는 원리는
언어 기반 보상 설계, 행동 조언, 계층적 계획 시스템에도 적용 가능성이 있다.
